% 张三慧V2第23章思考题23.1--23.9独立候选解，图23.30仍需TikZ复核。
\section*{第23章\quad 光的衍射：思考题解析}
\paragraph{思考题 23.1}
衍射是否显著取决于障碍物或孔口尺度 \(a\)
与波长 \(\lambda\) 的比值。
声波在空气中波长可达厘米至数米，
常与门窗、墙边或物体宽度相当，
因此容易绕过障碍和传播到几何阴影区。
可见光波长约 \(400\)--\(700\,\mathrm{nm}\)，
相对日常厘米—米障碍极小，
\[
\boxed{\lambda_{\rm sound}/a
\gg \lambda_{\rm visible}/a}.
\]
所以相同大小的障碍下声波衍射明显，
光的衍射角通常很小；光仍具有衍射本性。

\paragraph{思考题 23.2}
单缝远场各角度 \(\theta\) 的平面波分量，
经焦距 \(f\) 的理想会聚透镜后聚到
其后焦平面不同横坐标
\[
\boxed{x=f\tan\theta\simeq f\theta}.
\]
因此透镜把远场\emph{角度谱}转为可直接观察的
焦平面位置分布。
在理想夫琅禾费条件下
\(I(x)\propto\mathrm{sinc}^2(\pi a\sin\theta/\lambda)\)。
透镜不是制造衍射的根本原因：
衍射由有限狭缝上的相干波叠加产生。

\paragraph{思考题 23.3}
相邻半波带的平均相位相差约 \(\pi\)，
且带面积相同下振幅近似相等。
虽然第1带与第3带相差 \(2\pi\) 而同相，
第2带与第4带也同相，
但是\emph{第1/3带与第2/4带}
整体相差 \(\pi\)。
因此
\[
\boxed{E_1+E_2+E_3+E_4
\simeq A-A+A-A=0}.
\]
观察点P为单缝衍射暗纹；
不能只比较相隔一个带的两个分量，
忘了它们与另外一组的抵消。
精确分带积分须按原图23.30的斜光程进行。

\paragraph{思考题 23.4}
(1) 当\emph{平面波入射方向不变}、
透镜和观察屏固定时，把狭缝沿横向移动
\(x_0\) 只让远场振幅增加线性相位因子
\[
\tilde E(k_x)\mapsto
\tilde E(k_x)e^{-ik_xx_0},
\]
从而
\(\boxed{I(k_x)=|\tilde E(k_x)|^2\text{ 不变}}\)。
(2) 若横移\emph{有限距离的光源}
使入射波倾角改变为 \(\theta_i\)，
场因子包含入射线性相位，
衍射图样相对轴线整体偏转，
中央最大约满足
\(\theta\simeq\theta_i\)。
\[
\boxed{\text{平移缝本身不改变角谱；
改变入射角会平移焦面图样。}}
\]
如果移动狭缝同时显著改变照明均匀性，
则还需考虑非理想入射光束。

\paragraph{思考题 23.5}
中央明纹两侧第一暗纹满足
\(a\sin\theta_1=\pm\lambda\)，
小角时中央角全宽
\[
\boxed{\Delta\theta\simeq2\lambda/a}.
\]
随着缝宽 \(a\) 增大，衍射中央主瓣变窄，
\emph{角分辨能力提高}；
在相同入射照度下，总通过功率
\(\propto a\) 增加（狭缝高度固定）。
峰值强度、总通量和角宽不是同一物理量，
不能推断图样所有方向均同时变亮。

\paragraph{思考题 23.6}
若波长改为
\(\lambda\sim1\,\mathrm{cm}\)，
而眼瞳口径仍 \(D\sim3\,\mathrm{mm}\)，
则
\[
\boxed{\lambda/D\sim3.3>1}.
\]
传统圆孔Airy小角瑞利式
\(\theta_R\simeq1.22\lambda/D\)
已不适合解释为一个小于90°的分辨角；
物体不同方向的传播分量
无法被如此小的孔径有效区分。
因此\emph{无法像可见光一样通过瞳孔成清晰的
外部物体轮廓像}。
对小于波长的孔口还应使用完整的亚波长孔径
电磁理论及低耦合效率模型，
不能把形式值 \(\theta_R>1\) 当成真实的角度。

\paragraph{思考题 23.7}
相邻缝距 \(d\) 固定时，
理想主极大方向满足
\[
\boxed{d\sin\theta_m=m\lambda}
\]
只依赖 \(d,\lambda,m\)，
不因缝数 \(N\) 改变而移动。
有限N阵列因子
\(I\propto[\sin(N\beta)/\sin\beta]^2\)，
\(\beta=\pi d\sin\theta/\lambda\)，
使主峰角宽量级\(\Delta\theta\propto1/N\)
变窄，峰值若每缝振幅固定可\(\propto N^2\)。
单缝包络的\emph{缺级}由
\(a\sin\theta=p\lambda\) 与主峰条件重合
决定；若 \(a,d\) 均固定，
改变N不能消除几何缺级。
须区分光栅分辨本领 \(mN\) 与单缝宽度的角色。

\paragraph{思考题 23.8}
对相同狭缝的两条单缝复振幅
\(E_1,E_2\)，分别只开其中一条时：
\[
I_1(\theta)=|E_1|^2,\qquad
I_2(\theta)=|E_2|^2,
\]
两者各呈单缝包络
\(\mathrm{sinc}^2(\pi a\sin\theta/\lambda)\)。
两缝同时打开，
\[
\boxed{I_{12}=|E_1+E_2|^2
=I_1+I_2+2\Re(E_1E_2^*)}.
\]
当有稳定相位关系时，交叉项不为零，
在单缝包络内产生密集明暗干涉条纹。
若各缝之间随机失相干而长时间平均，
交叉项可消失，此时才有
\(\langle I_{12}\rangle
\simeq\langle I_1\rangle+\langle I_2\rangle\)。

\paragraph{思考题 23.9}
非周期缝的位置 \(x_j\) 可写为
\[
E(k_x)\propto A_{\rm slit}(k_x)
\sum_{j=1}^{N}e^{ik_xx_j},
\quad
I(k_x)\propto|E(k_x)|^2.
\]
对于\emph{等间距}缝，
各相位差在 \(k_xd=2\pi m\) 时共同重合，
形成锐利的非零级主极大。
若相邻缝距随机且无长程周期性，
通常\emph{不会有按固定级次规律排列的
清晰锐利光栅主峰}；
可能出现不规则的散斑或较平滑的衍射背景。
在\emph{正前方} \(k_x=0\)
所有缝位置相位仍相同，
\[
\boxed{E(0)\propto N,\qquad I(0)\propto N^2}
\]
（在单缝幅度与相干照明固定条件下），
因此中心相干最大通常仍在。
